\documentclass[12pt]{article}

%\usepackage[cm]{fullpage}
\setlength{\topmargin}{-1in}
\setlength{\oddsidemargin}{-0.25in}
\setlength{\evensidemargin}{-0.25in}
\setlength{\textwidth}{7in}
\setlength{\textheight}{10in}

\usepackage{array}
\usepackage{enumitem}
\usepackage{amssymb}
\usepackage{amsmath}
\usepackage{epstopdf}
\usepackage[T1]{fontenc}
\usepackage{graphics}
\usepackage{graphicx}
\usepackage{hyperref}
\usepackage{latexsym}
\usepackage{mathrsfs}
\usepackage{multicol}
\usepackage{setspace}
\usepackage{subfigure}
\usepackage{time}
\usepackage{times}

\newcommand{\To}{\Rightarrow}
\newcommand{\del}{\nabla}
\newcommand{\R}{\mathbb{R}}
\newcommand{\C}{\mathcal{C}}
\newcommand{\N}{\mathcal{N}}
\newcommand{\A}{\mathcal{A}}
\renewcommand{\O}{\mathcal{O}}
% \renewcommand{\N}{\mathbb{N}}
\newcommand{\ben}{\begin{enumerate}}
\newcommand{\een}{\end{enumerate}}
\newcommand{\eps}{\varepsilon}
\renewcommand{\epsilon}{\varepsilon}
\newcommand{\p}{\partial}
\newcommand{\<}{\langle}
\renewcommand{\>}{\rangle}
\renewcommand{\subset}{\subseteq}
\newcommand{\cont}{\Rightarrow \Leftarrow}
\newcommand{\back}{\backslash}
\newcommand{\norm}[1]{\|{#1}\|}
\newcommand{\abs}[1]{|{#1}|}
\newcommand{\ip}[1]{\langle{#1}\rangle}
\newcommand{\m}[1]{{\mathbf{#1}}}
\newcommand{\sts}{\m{S}^\textrm{T}\m{S}}
\newcommand{\rank}{\textrm{rank}}
\newcommand{\X}{\mathcal{X}}
\newcommand{\Y}{\mathcal{Y}}
\def\ds{\displaystyle}
\def\d{\partial}
\renewcommand{\arraystretch}{1.5}
\renewcommand{\L}{\mathscr{L}}

%\includegraphics[scale = 0.65, trim = 1in 3.25in 1in 3in]{blank_axes.pdf}


\pagestyle{empty}
\begin{document}
\noindent\rule{7in}{0.4pt}\\
\begin{center}
{\sc Differential Equations\\
Family of Solutions}
\end{center}
\noindent\rule{7in}{0.4pt}
%%%%%%%%%%%%%%%%%%%%%%%%%%%%%%%%%%%%%%%%%%%%%%%%
\section*{Family of solutions, number of constants}
\doublespacing

\noindent  We can actually use Example 3 in the previous section to start developing a couple of 
other definitions.  Recall that we showed that for any choice of constants $c_1$ and $c_2$, the function
$y = c_1e^{-3t} + c_2e^{4t}$ is a solution to the DE $y'' - y' - 12y = 0.$  Suppose we choose $c_1 = 2$
and $c_2 = 0.$  Then we have $y = 2e^{-3t}$.  In Example 3 we showed that this was a solution, but 
it turns out this was just a specific case of the solution $y = c_1e^{-3t} + c_2e^{4t}$.  Can you see that 
the other solution in Example 3, $y = e^{4t}$, is also just a specific case of the solution 
$y = c_1e^{-3t} + c_2e^{4t}$?\\

\noindent  It turns out that, as demonstrated in Example 3, \\ %DEs generally have an infinite number of 
%solutions.  
\line(1,0){500}. \\ When we solve DEs, we want to know about all of these solutions.\\

\noindent In the context of Example 3, the function $y = c_1e^{-3t} + c_2e^{4t}$ represents the 
%{\bf family of solutions} 
\line(1,0){200}\\ 
of the DE $y'' - y' - 12y = 0.$  Another term we use for the family of solutions is the 
%{\bf general solution}.  
\line(1,0){160}. \\
That means that the following are equivalent ways to ask exactly the same
question:
	\begin{itemize}
	\item  \mbox{}\\
	%Solve the DE $y'' - y' - 12y = 0.$
	\item  \mbox{}\\
	%Find the family of solutions for $y'' - y' - 12y = 0.$
	\item  \mbox{}\\
	%Find the general solution to $y'' - y' - 12y = 0.$
	\end{itemize}
\noindent  While we have not yet discussed {\it how} to find the solution, the correct answer to
all three of these would be $y = c_1e^{-3t} + c_2e^{4t}$ (as we demonstrated in Example 3).
%Margin note??  Note that  the ``$y=$'' is not optional.  The solution to a DE should be a function,
%and we need to include the dependent variable to make it clear to the reader that this is a
%function and not an expression.
\\

\noindent  Notice in Example 3 that the DE $y'' - y' - 12y = 0$ is a second order DE, and that the
general solution contains two constants.  This is not a coincidence; the order of the DE 
typically dictates how many generic constants are in the solution.  \line(1,0){375}\\ \\
\line(1,0){500}\\ \\
\line(1,0){500}\\
%The solution to a first order DE will generally have one 
%constant; the solution to a second order DE will have two constants; etc.\\


\end{document}